% Image credits for the photographs and AI illustrations of Book 5
% (University Biology, Year 3). Machine-readable license records live in
% images/book5/CREDITS.md; keep the two files in sync.
\chapter*{चित्र-श्रेय}

% Under bidi=basic a Latin run is ONE directional box, so a long author or
% institution name cannot be broken at all; \allowbreak inside it does nothing.
% The Arabic edition reported this page as the book's only overfull boxes.
% \sloppy must be in force when the paragraph is BROKEN, hence the \par before
% \endgroup -- without it the group closes first and the tolerance is restored.
\begingroup\sloppy

इस पुस्तक के रेखाचित्र मौलिक हैं। छायाचित्र अपने-अपने मुक्त अनुज्ञापत्रों के
अंतर्गत प्रयुक्त हैं; उनके रचयिताओं के प्रति आभार:

\begin{itemize}
  \item अध्याय~\ref{ch:b3:genomics}: Frederick Sanger का छायाचित्र ---
        सार्वजनिक अधिकार-क्षेत्र (Wikimedia Commons)।
  \item अध्याय~\ref{ch:b3:genetic-engineering}: साधारण चावल के साथ
        स्वर्णिम चावल, International Rice Research Institute, 2011 ---
        CC~BY~2.0 (Wikimedia Commons)।
  \item अध्याय~\ref{ch:b3:structural-biology}: Max Perutz का छायाचित्र,
        Associated Press, 1962 --- सार्वजनिक अधिकार-क्षेत्र
        (Wikimedia Commons)।
  \item अध्याय~\ref{ch:b3:bacteriology}: Robert Koch का चित्र, Wellcome
        Collection --- CC~BY~4.0 (Wikimedia Commons); अपनी प्रयोगशाला
        में Alexander Fleming का छायाचित्र, U.S. Navy Bureau of Medicine
        and Surgery का अभिलेखागार --- कोई ज्ञात अधिकार-प्रतिबंध नहीं
        (Wikimedia Commons)।
  \item अध्याय~\ref{ch:b3:innate-immunity}: Élie Metchnikoff का
        छायाचित्र, Bain News Service, Library of Congress ---
        सार्वजनिक अधिकार-क्षेत्र (Wikimedia Commons)।
  \item अध्याय~\ref{ch:b3:adaptive-immunity}: Edward Jenner का उत्कीर्ण
        चित्र, 1807 --- सार्वजनिक अधिकार-क्षेत्र (Wikimedia Commons)।
  \item अध्याय~\ref{ch:b3:nervous-systems}: Santiago Ramón y Cajal
        द्वारा बनाया पुरकिंजे कोशिका का रेखाचित्र, 1899 --- सार्वजनिक
        अधिकार-क्षेत्र (Wikimedia Commons)।
  \item अध्याय~\ref{ch:b3:endocrinology}: Frederick Banting और Charles
        Best का छायाचित्र, लगभग 1924 --- सार्वजनिक अधिकार-क्षेत्र
        (Wikimedia Commons)।
  \item अध्याय~\ref{ch:b3:molecular-evolution}: Motoo Kimura का
        छायाचित्र, 1986, छायाकार अज्ञात, S.~Kumar एवं T.~Gojobori,
        \emph{Molecular Biology and Evolution} (2023) के माध्यम से ---
        CC~BY~4.0 (Wikimedia Commons)।
  \item अध्याय~\ref{ch:b3:behavioural-ecology}: Rob Mieremet द्वारा
        लिया Niko Tinbergen का छायाचित्र, Anefo, 1973 --- CC0,
        Nationaal Archief (Wikimedia Commons)।
  \item अध्याय~\ref{ch:b3:conservation-biology}: अंतिम प्रवासी कबूतर
        मार्था का छायाचित्र, Enno Meyer, 1912 --- सार्वजनिक
        अधिकार-क्षेत्र (Wikimedia Commons); Yellowstone National Park
        में भेड़िये, National Park Service, 1996 --- सार्वजनिक
        अधिकार-क्षेत्र (Wikimedia Commons); काकापो सिरोक्को, Department
        of Conservation, न्यूज़ीलैंड, 2012 --- CC~BY~2.0
        (Wikimedia Commons)।
\end{itemize}

प्रत्येक छायाचित्र के पूर्ण स्रोत-लिंक और अनुज्ञापत्र के पते पुस्तक के
भंडार में \texttt{images/\allowbreak book5/\allowbreak CREDITS.md} में दिए गए हैं।

\bigskip
छायाचित्रों और मौलिक रेखाचित्रों के साथ-साथ, इस खंड के अध्यायों में फैले
लगभग नब्बे चित्र कृत्रिम बुद्धिमत्ता से बनी चित्रांकन-कृतियाँ हैं, जो
पुस्तक के संपादकों द्वारा लिखे संकेतों से OpenAI की चित्र-जनन प्रणाली से
बनाई गईं और सम्मिलित करने से पहले जैविक शुद्धता के लिए जाँची गईं।
बनाए गए हर चित्र का पूरा संकेत इस किताब के भंडार में, फ़ाइल
\texttt{images/\allowbreak book5/\allowbreak ai/\allowbreak PROMPTS.md} में दर्ज है।
\par\endgroup
\clearpage
